\chapter{ทฤษฎีที่เกี่ยวข้อง}
\label{Ch:Background}
การศึกษานี้ตั้งอยู่บนพื้นฐานของแนวคิดด้านพลังงานของซอฟต์แวร์ (Software Energy Consumption) แนวคิดด้าน Green Software Engineering และงานวิจัยหลายชิ้นที่สำรวจผลกระทบของ Design Patterns, Code Smells และ Refactoring ต่อประสิทธิภาพด้านพลังงานของระบบ เพื่อพัฒนาไลบรารีสำหรับตรวจจับ Green Code Smell ในภาษา Python อย่างเป็นระบบ งานวิจัยที่เกี่ยวข้องสามารถสรุปได้ดังนี้

% -----------------------------------------------------

\section{Green Software Engineering และพลังงานของซอฟต์แวร์}
งานสำรวจเชิงระบบโดย Poy et al.~\cite{poy2025impact} แสดงให้เห็นว่าการเลือกใช้ design patterns, การเกิด code smells และการ refactoring มีความสัมพันธ์โดยตรงกับการใช้พลังงานของซอฟต์แวร์ โดยพบว่าบาง design patterns สามารถช่วยลดการใช้พลังงานได้ในขณะที่บาง patterns อาจเพิ่มพลังงานขึ้นหากใช้งานไม่เหมาะสมซึ่งชี้ให้เห็นถึงความสำคัญของการออกแบบโค้ดให้มีความยั่งยืนตั้งแต่ต้นทาง

นอกจากนี้ Le Goaër~\cite{leGoaer2024tactics} ได้เสนอแนวคิด Green Tactics, Green Patterns และ Green Smells เพื่อใช้เป็นหลักการในการลดการทำงานที่สิ้นเปลืองทรัพยากรและเป็นรากฐานสำคัญของแนวทาง Green Code Smell Detection ในงานนี้

% -----------------------------------------------------

\section{Code Smells และผลกระทบด้านพลังงาน}
งานของ Verdecchia et al.~\cite{verdecchia2018energy} ถือเป็นหลักฐานเชิงประจักษ์สำคัญที่แสดงว่าการ refactor code smells สามารถลดการใช้พลังงานได้อย่างมีนัยสำคัญ โดยเฉพาะ code smells ประเภท Feature Envy และ Long Method ซึ่งเมื่อนำไป refactor แล้วให้ผลลัพธ์ลดการใช้พลังงานได้มากถึง 49.9\% ผลลัพธ์ดังกล่าวสนับสนุนสมมติฐานว่าการปรับโครงสร้างโค้ด สามารถส่งผลโดยตรงต่อพฤติกรรมด้านพลังงานของโปรแกรม

อย่างไรก็ตาม งานวิจัยดังกล่าวยังชี้ให้เห็นว่าไม่ใช่ทุก code smell ที่ให้ผลด้านพลังงานในทิศทางเดียวกันเช่น การ refactor God Class อาจเพิ่มจำนวนข้อความส่งระหว่างออบเจกต์ นำไปสู่การใช้พลังงานเพิ่มขึ้นแสดงให้เห็นถึงความซับซ้อนของความสัมพันธ์ระหว่าง “คุณภาพโค้ด” และ “ประสิทธิภาพด้านพลังงาน”

% -----------------------------------------------------

\section{Green Code Smells และ ecoCode}
ระบบ ecoCode ของ Le Goaër~\cite{legoaer2023ecocode} ได้ขยายความสามารถของ SonarQube เพื่อรองรับการตรวจจับ \textit{Green Code Smells} ซึ่งถูกนิยามว่าเป็น “รูปแบบการเขียนโค้ดที่ส่งผลต่อการใช้พลังงานหรือเพิ่ม Carbon Footprint ของซอฟต์แวร์” เครื่องมือนี้เน้นการตรวจจับปัญหาที่เกี่ยวข้องกับการใช้ API, ทรัพยากรฮาร์ดแวร์ และรูปแบบการประมวลผลที่ไม่เหมาะสมซึ่งถือเป็นฐานข้อมูลสำคัญในการประยุกต์ใช้กับภาษา Python ในงานนี้

เวอร์ชันที่ถูกเก็บใน HAL Archive~\cite{legoaer2023hal} ยังได้อธิบายโครงสร้างของ Green Code Smell Catalogue ซึ่งเป็นชุดกฎที่ไม่ได้ผูกติดกับภาษาใดภาษาเดียว แต่เน้นการวิเคราะห์โครงสร้างโค้ดในเชิงสถิติและพฤติกรรมการใช้พลังงาน ด้วยเหตุนี้ แนวคิดของ ecoCode จึงถูกนำมาใช้เป็นแม่แบบสำหรับการออกแบบกฎตรวจจับ Green Code Smell ของโปรเจกต์นี้

% -----------------------------------------------------

\section{มาตรวัดและการประเมินผลประสิทธิภาพการปล่อยคาร์บอน (Metrics and Evaluation)}
การประเมินความยั่งยืนของซอฟต์แวร์จำเป็นต้องมีดัชนีชี้วัดที่ชัดเจน เพื่อเปลี่ยนจากแนวคิดเชิงคุณภาพให้เป็นเชิงปริมาณ โดยมาตรฐานและแนวคิดที่สำคัญที่เกี่ยวข้องกับโครงงานนี้มีดังต่อไปนี้

\subsection{ข้อกำหนด Software Carbon Intensity (SCI)}
{\boldfont Software Carbon Intensity (SCI)} เป็นมาตรฐานที่นำเสนอโดย Green Software Foundation เพื่อใช้ในการคำนวณอัตราการปล่อยก๊าซเรือนกระจกของซอฟต์แวร์ต่อหนึ่งหน่วยการทำงาน (Functional Unit) โดยดัชนี SCI เน้นที่ ``ความเข้มข้น'' ของการปล่อยคาร์บอน และสามารถนำเสนอได้ในรูปของสมการต่อไปนี้~\cite{sciSpecification}

\[
SCI = \frac{(E \times I) + M}{R}
\]

คำอธิบายตัวแปรมีดังนี้
\begin{itemize}
    \item {\boldfont $E$ (Energy)}: พลังงานไฟฟ้าที่ซอฟต์แวร์ใช้ หน่วยเป็น kilowatt-hours (kWh)
    \item {\boldfont $I$ (Location-based Marginal Carbon Intensity)}: ค่าความเข้มข้นของคาร์บอนในสายส่งไฟฟ้า ณ เวลาและสถานที่นั้น ๆ หน่วยเป็น $gCO_2e/kWh$
    \item {\boldfont $M$ (Embodied Emissions)}: คาร์บอนที่เกิดขึ้นจากการผลิตและการจัดการวงจรชีวิตของฮาร์ดแวร์
    \item {\boldfont $R$ (Functional Unit)}: หน่วยวัดการทำงานของซอฟต์แวร์ เช่น ต่อจำนวนผู้ใช้งาน (User), ต่อการเรียกใช้ API (API Call) หรือต่อระยะเวลา (Time)
\end{itemize}

แนวคิดของ SCI มีความสำคัญต่อโครงงานนี้ เนื่องจากช่วยให้การประเมินผลด้านความยั่งยืนของซอฟต์แวร์ไม่หยุดอยู่เพียงการวัดพลังงานที่ใช้ทั้งหมด แต่สามารถเปรียบเทียบการปล่อยคาร์บอนต่อผลลัพธ์ของการทำงานได้อย่างเป็นระบบและเป็นธรรมมากขึ้น

\subsection{แนวคิดและวิธีการของ CodeCarbon}
CodeCarbon เป็นไลบรารีประเภท open-source สำหรับภาษา Python ที่ออกแบบมาเพื่อติดตามและประมาณการปริมาณการปล่อยก๊าซคาร์บอนไดออกไซด์เทียบเท่า ($CO_2e$) จากการประมวลผลแบบ real-time โดยมีหลักการคำนวณเบื้องต้นดังนี้~\cite{codecarbonMethodology}

\[
\text{Carbon Emissions} = \text{Energy Consumed} \times \text{Carbon Intensity}
\]

กลไกการทำงานสำคัญของ CodeCarbon สามารถสรุปได้ดังนี้
\begin{itemize}
    \item {\boldfont การติดตามการใช้พลังงาน (Energy Tracking)}: ดึงข้อมูลการใช้พลังงานจากฮาร์ดแวร์ผ่านอินเทอร์เฟซที่เกี่ยวข้อง เช่น Intel Power Gadget หรือ Running Average Power Limit (RAPL) สำหรับ CPU และ NVIDIA Management Library (NVML) สำหรับ GPU
    \item {\boldfont การคำนวณตามพื้นที่ทางภูมิศาสตร์ (Localization)}: ระบุตำแหน่งของผู้ใช้ผ่าน IP address เพื่อนำไปอ้างอิงกับฐานข้อมูลความเข้มข้นของคาร์บอนในโครงข่ายไฟฟ้าของแต่ละพื้นที่
    \item {\boldfont การแปลงเป็นค่า $CO_2$ Equivalent ($CO_2e$)}: นำค่าพลังงานที่วัดได้คูณกับค่าความเข้มข้นของคาร์บอนเพื่อประเมินออกมาเป็นปริมาณก๊าซเรือนกระจกที่ปล่อยออกมาจริง
\end{itemize}

ในบริบทของโครงงานนี้ CodeCarbon ทำหน้าที่เป็นเครื่องมือสำหรับสนับสนุนการประเมินผลเชิงปริมาณ โดยเชื่อมโยงผลการตรวจจับ Green Code Smell เข้ากับข้อมูลการใช้พลังงานและการปล่อยคาร์บอนที่สามารถวัดและเปรียบเทียบได้

\subsection{มาตรวัดขนาดโปรเจกต์ (Project Size Measurement)}
การกำหนดหน่วยวัดขนาดของซอฟต์แวร์ถือเป็นปัจจัยสำคัญในการทำ benchmarking และใช้เป็นค่า $R$ (Functional Unit) ในสมการ SCI โดยมาตรวัดที่นำมาพิจารณาในโครงงานนี้มี 3 รูปแบบหลัก ได้แก่~\cite{vanHeeringen2015size}
\begin{itemize}
    \item {\boldfont Lines of Code (LOC)}: การวัดเชิงปริมาณด้วยการนับจำนวนบรรทัดของโค้ด เป็นวิธีที่ง่ายและตรงไปตรงมา แต่ไม่สามารถสะท้อนคุณภาพหรือความซับซ้อนของงานที่ซอฟต์แวร์ทำจริงได้อย่างครบถ้วน
    \item {\boldfont Per Execution}: การวัดเชิงพฤติกรรมโดยนับตามรอบการประมวลผล เหมาะสำหรับการทดสอบประสิทธิภาพในระดับ \textit{micro-benchmark} ของฟังก์ชันหรือกระบวนการเฉพาะจุด
    \item {\boldfont Cosmic Function Points (CFP)}: วิธีการวัดขนาดซอฟต์แวร์เชิงฟังก์ชันตามมาตรฐาน ISO/IEC 19761 ซึ่งพัฒนาโดย COSMIC โดยมีแนวคิดว่าขนาดของซอฟต์แวร์สามารถวัดได้จากจำนวนครั้งที่ข้อมูลถูกเคลื่อนย้ายภายใน \textit{Functional Process} หรือหน่วยของงานที่ผู้ใช้ต้องการ~\cite{abran2018cosmic}
\end{itemize}

สำหรับ CFP การเคลื่อนย้ายข้อมูลถูกแบ่งออกเป็น 4 ประเภทหลัก ได้แก่ Entry (E), Exit (X), Read (R) และ Write (W) โดยการเคลื่อนย้ายข้อมูลแต่ละครั้งมีค่าเท่ากับ 1 CFP เสมอ ทำให้การวัดมีความเป็นกลาง ไม่ขึ้นกับภาษาโปรแกรมหรือเทคโนโลยี และสามารถประเมินได้ตั้งแต่ช่วงวิเคราะห์ความต้องการก่อนเริ่มพัฒนา

\begin{table}[H]
\centering

\begin{adjustbox}{max width=\textwidth}
\renewcommand{\arraystretch}{1.3}
\setlength{\tabcolsep}{6pt}

\begin{tabular}{|p{2.8cm}|p{3.8cm}|p{4.8cm}|p{4.8cm}|}
\hline
{\boldfont ประเภท} &
{\boldfont ทิศทาง} &
{\boldfont ความหมาย} &
{\boldfont ตัวอย่าง} \\
\hline

E - Entry
&
User $\rightarrow$ Process
&
ข้อมูลที่ผู้ใช้หรือระบบภายนอกส่งเข้ามาให้ Process ประมวลผล
&
กรอกฟอร์มลงทะเบียน, ส่ง JSON จาก API
\\
\hline

X - Exit
&
Process $\rightarrow$ User
&
ข้อมูลที่ Process ส่งออกไปให้ผู้ใช้หรือระบบภายนอก
&
แสดงผลลัพธ์บนหน้าจอ, ส่ง response กลับ
\\
\hline

R - Read
&
Storage $\rightarrow$ Process
&
Process ดึงข้อมูลจาก Persistent Storage มาใช้งาน
&
ค้นหาข้อมูลลูกค้าจากฐานข้อมูล
\\
\hline

W - Write
&
Process $\rightarrow$ Storage
&
Process บันทึกข้อมูลลง Persistent Storage
&
INSERT/UPDATE ลงฐานข้อมูล, เขียนลงไฟล์
\\
\hline

\end{tabular}
\end{adjustbox}

\captionsetup{justification=centering, name=ตาราง}
\caption{การเคลื่อนย้ายข้อมูล 4 ประเภทของ COSMIC Function Points (CFP)}
\label{table:cfp-data-movement}
\end{table}

เมื่อเปรียบเทียบมาตรวัดทั้งสามรูปแบบ LOC เหมาะกับการวัดเชิงโครงสร้างในระดับพื้นฐาน, Per Execution เหมาะกับการวัดผลในเชิงพฤติกรรมการทำงานของโปรแกรม ขณะที่ CFP เหมาะสำหรับใช้เป็นหน่วยวัดเชิงฟังก์ชันเพื่อรองรับการประเมินตามกรอบ SCI ได้อย่างเป็นระบบมากกว่า เนื่องจากสะท้อน ``งานที่ซอฟต์แวร์ส่งมอบ'' มากกว่าปริมาณโค้ดเพียงอย่างเดียว~\cite{vanHeeringen2015size,abran2018cosmic}

% -----------------------------------------------------

\section{สรุป}
จากงานวิจัยและแนวคิดที่เกี่ยวข้อง สามารถสรุปประเด็นสำคัญได้ดังนี้

\begin{itemize}
    \item Code Smells มีผลกระทบต่อพลังงานอย่างมีนัยสำคัญ โดยเฉพาะในกรณีของ Feature Envy และ Long Method~\cite{verdecchia2018energy}
    \item การใช้ design patterns และโครงสร้างโค้ดที่เหมาะสมมีส่วนช่วยลดพลังงานได้~\cite{poy2025impact}
    \item แนวคิด Green Tactics และ Green Patterns เป็นพื้นฐานสำคัญของ Green Code Smell Rules~\cite{leGoaer2024tactics}
    \item ecoCode เป็นตัวอย่างของ static analyzer ที่ตรวจจับ Green Coding ได้จริงในระบบอุตสาหกรรม~\cite{legoaer2023ecocode}
    \item แนวคิด Software Carbon Intensity (SCI) ช่วยให้การประเมินความยั่งยืนของซอฟต์แวร์สามารถทำได้ในเชิงปริมาณ โดยพิจารณาการปล่อยคาร์บอนต่อหนึ่งหน่วยการทำงาน ทำให้สามารถเปรียบเทียบประสิทธิภาพด้านสิ่งแวดล้อมของซอฟต์แวร์ได้อย่างเป็นระบบและเป็นรูปธรรมมากขึ้น~\cite{sciSpecification}
    \item CodeCarbon เป็นเครื่องมือสำคัญที่ช่วยสนับสนุนการวัดค่าการใช้พลังงานและการปล่อยก๊าซเรือนกระจกจากการประมวลผลจริง ขณะที่แนวคิดด้านการวัดขนาดซอฟต์แวร์ เช่น LOC, Per Execution และโดยเฉพาะ Cosmic Function Points (CFP) มีบทบาทสำคัญในการกำหนดหน่วยการทำงานสำหรับการประเมินตามกรอบ SCI~\cite{codecarbonMethodology,vanHeeringen2015size,abran2018cosmic}
\end{itemize}

องค์ความรู้ดังกล่าวถูกนำมาประยุกต์เพื่อออกแบบชุดกฎสำหรับการตรวจจับ Green Code Smell ในภาษา Python โดยใช้เทคนิค Abstract Syntax Tree (AST) ควบคู่กับแนวทางการประเมินผลด้านพลังงานและคาร์บอน เพื่อสร้างเครื่องมือที่ช่วยผู้พัฒนาในการเขียนโค้ดที่ประหยัดพลังงาน ยั่งยืน และสามารถประเมินผลกระทบต่อสิ่งแวดล้อมได้อย่างเป็นรูปธรรม